% Generated by build_report.py. Edit the shared content there, then regenerate.
\documentclass[10pt,a4paper]{article}
\usepackage{fontspec}
\setmainfont{texgyretermes}[Extension=.otf,UprightFont=*-regular,BoldFont=*-bold,ItalicFont=*-italic,BoldItalicFont=*-bolditalic]
\usepackage[margin=1.9cm,headheight=14pt]{geometry}
\usepackage{amsmath,amssymb,booktabs,tabularx,longtable,array,ragged2e,enumitem,microtype,xcolor,tikz,fancyhdr,graphicx}
\usetikzlibrary{arrows.meta,positioning}
\usepackage[hidelinks,unicode]{hyperref}
\definecolor{ink}{HTML}{183A45}\definecolor{teal}{HTML}{087F7A}\definecolor{pale}{HTML}{EAF4F2}\definecolor{gray}{HTML}{56656B}
\newcolumntype{P}[1]{>{\RaggedRight\arraybackslash}p{#1}}
\setlength{\parindent}{0pt}\setlength{\parskip}{5pt}\setlength{\emergencystretch}{2em}
\setlength{\tabcolsep}{4pt}\renewcommand{\arraystretch}{1.16}
\setlist{nosep,leftmargin=1.4em,topsep=4pt}
\pagestyle{fancy}\fancyhf{}\fancyhead[L]{\small\color{gray}BẢO VỆ RIÊNG TƯ QUỸ ĐẠO}\fancyhead[R]{\small\color{gray}03/10/2026}\fancyfoot[C]{\small\thepage}\renewcommand{\headrulewidth}{0.3pt}
\newcommand{\key}[1]{\par\smallskip\noindent\colorbox{pale}{\parbox{\dimexpr\linewidth-2\fboxsep}{\textbf{#1}}}\par\smallskip}
\hypersetup{pdftitle={Bảo vệ riêng tư quỹ đạo: dữ liệu, phương pháp và kết quả},pdfauthor={}}
\begin{document}
\begin{center}{\LARGE\bfseries\color{ink}Bảo vệ riêng tư quỹ đạo}\\[5pt]{\large Related works, kiến trúc và kết quả đối chứng}\\[4pt]{\small Bản trình bày rút gọn · 03/10/2026 · Geo-I là phương pháp chính · Hai phiên bản benchmark báo riêng}\end{center}
\section{Related works: mức bao phủ S1–S10}\label{sec:1}
Khảo sát giữ cửa sổ \textbf{21/09/2023–21/09/2026}. Bảng đối chiếu từng paper với mười nhiệm vụ của ta; mức bao phủ phải đọc cùng giả định và giới hạn ở cột cuối.

S1: vị trí hiện tại; S2: nơi dừng; S3: đoạn đường đã đi; S4: liên kết danh tính giữa phiên; S5: cạnh đường kế tiếp; S6: đích chưa tới; S7: nội dung/ý định truy vấn; S8: định vị nhờ dữ liệu người đi cùng; S9: điểm đầu bị che; S10: điểm cuối bị che của chuyến đã hoàn tất.

\textbf{Fully (F):} paper trực tiếp xử lý mục tiêu suy luận của scenario, trong giả định của paper. \textbf{Partially (P):} chỉ xử lý một phần hoặc bối cảnh hẹp hơn. \textbf{CX:} chưa đủ nguồn để phân loại. \textbf{—:} chưa tìm thấy bằng chứng xử lý trực tiếp; không kết luận cơ chế chắc chắn thất bại. F/P là đối chiếu của ta, không phải nhãn do tác giả công bố, không có nghĩa bảo vệ 100\% hoặc đã vượt các ca A/B/C của ta.

\begingroup\small
\begin{longtable}{@{}P{3.387cm}P{1.742cm}P{2.613cm}P{8.614cm}@{}}
\toprule
\textbf{Paper} & \textbf{Fully} & \textbf{Partially / CX} & \textbf{Căn cứ và giới hạn}\\\midrule
\endfirsthead
\toprule
\textbf{Paper} & \textbf{Fully} & \textbf{Partially / CX} & \textbf{Căn cứ và giới hạn}\\\midrule
\endhead
\bottomrule\endfoot
TransProtect 2024 [\hyperlink{ref-R1}{R1}] & S1, S3 & S2 & §3–5: suy vị trí từ chuỗi đã làm nhiễu, biết đường/traffic; chưa đánh giá riêng nơi dừng.\\
\addlinespace[3pt]
Semantic correlation 2026 [\hyperlink{ref-R2}{R2}] & S1, S3 & S2 & §3–6: giữ dummy hợp lý theo thời gian/ngữ nghĩa; chưa chấm riêng nơi dừng hay ý định S7.\\
\addlinespace[3pt]
Fake queries 2026 [\hyperlink{ref-R3}{R3}] & S1, S3 & S2 & §3–6: chèn truy vấn giả để khó nối tuyến; chưa đánh giá riêng tọa độ nơi dừng.\\
\addlinespace[3pt]
Road-aware PTPPM 2025 [\hyperlink{ref-R4}{R4}] & S1, S3 & S2 & §III, VI: chống suy vị trí/quỹ đạo có tương quan; bước dự báo chưa là phép thử S5/S6.\\
\addlinespace[3pt]
CPCROK 2025 [\hyperlink{ref-R5}{R5}] & — & S4 & §3–5: chống nối bí danh xe trong VANET; chưa tách danh tính người/thiết bị như S4.A/B/C.\\
\addlinespace[3pt]
DP-FETC 2025 [\hyperlink{ref-R6}{R6}] & — & S3, S8, S9, S10 & SynFE/PreOD: sinh dữ liệu offline, đổi đầu/cuối của tuyến tương quan; không chấm định vị từ người đi cùng.\\
\addlinespace[3pt]
Improved PIR 2025 [\hyperlink{ref-R7}{R7}] & CX & CX: S1, S7 & Tóm tắt nêu bảo vệ vị trí/nội dung bằng mã hóa; thiếu toàn văn để đối chiếu quyền quan sát và suy ý định.\\
\addlinespace[3pt]
SoK 2024 [\hyperlink{ref-R8}{R8}] & Không áp dụng & Không áp dụng & Khảo sát cách đánh giá/tấn công; không phải cơ chế bảo vệ để gán F/P.\\
\addlinespace[3pt]
PRISM 2025 [\hyperlink{ref-R9}{R9}] & CX & CX: S1, S3 & Tóm tắt: bảo vệ chia sẻ quỹ đạo; chưa đủ mô hình attacker. Dùng LSTM không xác nhận xử lý S5/S6.\\
\addlinespace[3pt]
Options to Action 2026 [\hyperlink{ref-R10}{R10}] & Không áp dụng & Không áp dụng & Khảo sát sử dụng tính năng, có che đầu/cuối (S9/S10); không đo khả năng chống suy luận.\\
\addlinespace[3pt]
Forsch et al. 2023 [\hyperlink{ref-R11}{R11}] & S9, S10 & — & §3.2: S-TT cắt cả hai đầu, xét nhiều tuyến cùng đích; bảo đảm theo tập tấn công/địa điểm cho trước, offline.\\
\addlinespace[3pt]
EV querying 2024 [\hyperlink{ref-R12}{R12}] & S1, S3 & S2 & Proposed Query Model: AGeoI + dummy chống định vị/nối tuyến online; chưa chấm riêng nơi dừng hoặc S4–S10.\\
\end{longtable}
\endgroup
Scenario không được liệt kê ở một hàng được hiểu là —; với R7/R9 là chưa xác minh, với R8/R10 là không áp dụng. Bảng cho thấy các cơ chế chuỗi chủ yếu xét S1/S3; S2 cần phép thử nơi dừng riêng, còn S9/S10 cần kiểm tra suy ngược từ phần tuyến vẫn công bố.


\clearpage
\subsection{Metrics gốc và điều rút ra từ related works}\label{sub:1.1}
Bảng dưới ghi các chỉ số đã xác minh trong nguồn. $\uparrow$/$\downarrow$ là hướng tốt hơn theo mục tiêu của paper; CX nghĩa là chưa xác minh đủ định nghĩa. Các paper dùng dataset và nhiệm vụ khác nhau, nên không so trực tiếp điểm số gốc của chúng.

\begingroup\small
\begin{longtable}{@{}P{3.871cm}P{4.355cm}P{3.871cm}P{4.258cm}@{}}
\toprule
\textbf{Paper / hướng nghiên cứu} & \textbf{Privacy / chẩn đoán} & \textbf{Utility} & \textbf{Chi phí / phạm vi}\\\midrule
\endfirsthead
\toprule
\textbf{Paper / hướng nghiên cứu} & \textbf{Privacy / chẩn đoán} & \textbf{Utility} & \textbf{Chi phí / phạm vi}\\\midrule
\endhead
\bottomrule\endfoot
TransProtect 2024 [\hyperlink{ref-R1}{R1}] & EIE $\uparrow$: sai số suy luận & $\Delta$c $\downarrow$: chênh chi phí hành trình & Byte + độ trễ toàn dịch vụ: CX\\
\addlinespace[3pt]
Semantic correlation 2026 [\hyperlink{ref-R2}{R2}] & ASR $\uparrow$: đạt ẩn danh; DER $\uparrow$: dummy/ngữ nghĩa & Chưa đo top-5 POI bằng DER & Thời gian sinh $\downarrow$\\
\addlinespace[3pt]
Fake queries 2026 [\hyperlink{ref-R3}{R3}] & Số đường khó phân biệt $\uparrow$; ASR $\uparrow$ & Chất lượng POI: CX & Thời gian sinh, RSS $\downarrow$\\
\addlinespace[3pt]
Road-aware PTPPM 2025 [\hyperlink{ref-R4}{R4}] & Sai số $\uparrow$; tấn công đúng $\downarrow$ & Dịch chuyển đầu ra $\downarrow$ & Chi phí LBS đầu-cuối: CX; preprint\\
\addlinespace[3pt]
CPCROK 2025 [\hyperlink{ref-R5}{R5}] & Tỷ lệ khớp quỹ đạo $\rho$ $\downarrow$ & Bối cảnh VANET & Số quỹ đạo giả $\psi$ $\downarrow$; ước lượng gián tiếp chi phí\\
\addlinespace[3pt]
DP-FETC 2025 [\hyperlink{ref-R6}{R6}] & MI $\downarrow$; bảo đảm $\varepsilon$-DP & JSD phân bố $\downarrow$ & Độ phức tạp; xuất bản offline\\
\addlinespace[3pt]
Improved PIR 2025 [\hyperlink{ref-R7}{R7}] & Phân tích an toàn; metric tấn công: CX & Truy hồi cụ thể: CX & Chi phí tính toán; mới đủ tóm tắt\\
\addlinespace[3pt]
SoK 2024 [\hyperlink{ref-R8}{R8}] & Tấn công tái dựng/liên kết & Phân bố, hình học, tác vụ & Khảo sát, không phải model\\
\addlinespace[3pt]
PRISM 2025 [\hyperlink{ref-R9}{R9}] & LPI; công thức CX & TDD, DC; định nghĩa CX & Thời gian xử lý/phản hồi\\
\addlinespace[3pt]
Options to Action 2026 [\hyperlink{ref-R10}{R10}] & Tỷ lệ dùng tính năng & Nhận thức và sử dụng & Nghiên cứu hành vi\\
\addlinespace[3pt]
Forsch et al. 2023 [\hyperlink{ref-R11}{R11}] & Trình bày S-TT / k-site-unidentifiability & Che điểm nhạy cảm & Chương tổng quan; kế thừa S-TT 2022\\
\addlinespace[3pt]
EV querying 2024 [\hyperlink{ref-R12}{R12}] & AGeoI + dummy & Dịch vụ truy vấn trạm sạc & Chưa xác minh phép thử điểm đầu/cuối\\
\end{longtable}
\endgroup
\textbf{Mốc nền giữ riêng:} DLS 2014 [\hyperlink{ref-B1}{B1}], RDG 2021 [\hyperlink{ref-B2}{B2}]; AnotherMe [\hyperlink{ref-B3}{B3}] có số tạp chí 2024 nhưng online 11/09/2023, ngoài cửa sổ. S-TT 2022 [\hyperlink{ref-B4}{B4}], tấn công EPZ 2022 [\hyperlink{ref-B5}{B5}] và Data Stream Obfuscator 07/2023 [\hyperlink{ref-B6}{B6}] bổ sung góc nhìn điểm đầu/cuối. Không gọi các nguồn này là paper mới trong ba năm.

\key{Các paper đo những khía cạnh khác nhau. Để so trên bài toán của ta, cần cùng đo khả năng suy luận của đối thủ, chất lượng dịch vụ và chi phí.}


\clearpage
\subsection{Vì sao chọn bộ metrics đề xuất?}\label{sub:1.2}
Đầu ra khác nhau là một nguyên nhân: một vị trí thay thế, một quỹ đạo giả, tập K dummy hay danh sách query chèn. Metric dựa vào độ giống dummy hoặc entropy của một tập không tự áp dụng cho mọi đầu ra. Tuy nhiên, khác threat model, nhiệm vụ, dataset và mẫu số cũng làm điểm không tương đương. Ta chấm toàn bộ dữ liệu mà server/attacker được phép thấy, rồi yêu cầu suy cùng tọa độ đích; không chỉ chọn một dummy để chấm.

\begingroup\small
\begin{longtable}{@{}P{5.513cm}P{4.135cm}P{6.990cm}@{}}
\toprule
\textbf{Vấn đề từ related works} & \textbf{Metric đề xuất} & \textbf{Ý nghĩa trong bài toán của ta}\\\midrule
\endfirsthead
\toprule
\textbf{Vấn đề từ related works} & \textbf{Metric đề xuất} & \textbf{Ý nghĩa trong bài toán của ta}\\\midrule
\endhead
\bottomrule\endfoot
Độ bất định, tương đồng dummy và lượng thông tin rò rỉ đo các khía cạnh khác nhau. ASR cũng không có cùng định nghĩa giữa các paper. & Hit50/100/200 $\downarrow$; MAE, median, p90 $\uparrow$ & Đối thủ (attacker) phải đoán tọa độ thật từ dữ liệu công bố. Hit100 là tỷ lệ đoán cách đáp án $\le$100 m, dùng được cho cả tập dummy và vị trí thay thế.\\
\addlinespace[3pt]
MAE lớn vẫn có thể che một số lần suy đúng rất gần. & Đọc Hit cùng phân bố sai số & MAE là sai số trung bình; median là trung vị; p90 là ngưỡng chứa 90\% sai số mẫu. S1 chấm một điểm, S2 nơi dừng, S3 chuỗi điểm, S9/S10 điểm đầu/cuối.\\
\addlinespace[3pt]
Bảo toàn phân bố hoặc giảm độ lệch chưa cho biết người dùng có nhận đúng địa điểm cần tìm. & Recall@5 $\uparrow$; trả đủ $\uparrow$; $\Delta$D đường $\downarrow$ & So 5 địa điểm được chọn tại thiết bị với 5 địa điểm chuẩn theo GPS thật, cùng trạng thái máy chủ. Mỗi ca cần Recall $\ge$90\%.\\
\addlinespace[3pt]
Chỉ đếm dummy hoặc thời gian sinh sẽ bỏ sót phản hồi, truy vấn chèn và tải ngoài chuyến. & Byte yêu cầu + phản hồi $\downarrow$; latency p50/p95 $\downarrow$ & Đếm toàn bộ dữ liệu gửi/nhận cho mỗi lần dùng dịch vụ thật. Độ trễ toàn dịch vụ chưa đo; chưa có điểm tổng Q thực nghiệm.\\
\end{longtable}
\endgroup
POI là địa điểm như nhà thuốc hoặc trạm sạc. Recall@5 = số POI chuẩn tìm được / số POI chuẩn (tối đa 5). Ví dụ, tìm đúng 4/5 địa điểm được 80\%. Có đáp án mà không phục vụ thì tính 0; không có đáp án chuẩn thì để thiếu. “Trả đủ” đo số lượng; $\Delta$D đo khoảng cách đường tăng thêm.

Để so công bằng, giữ cùng mục tiêu, quyền quan sát và dịch vụ: máy chủ trả top-10, thiết bị chọn top-5. Với mỗi loại dữ liệu công bố (transcript), chọn bộ suy luận phù hợp trên nhóm phát triển rồi giữ cố định khi kiểm tra. Gộp các lần chạy trong bản ghi và nhóm tuyến; cho các ca trọng số bằng nhau trong scenario, rồi cho năm scenario trọng số bằng nhau. S10 chỉ có A/B. Chi phí truyền vẫn phải báo riêng.

\key{Bộ đo trả lời ba câu hỏi: đối thủ đoán đúng tới đâu, người dùng nhận đúng kết quả đến đâu và cần bao nhiêu chi phí. Đóng góp là cách chọn và áp dụng chung các chỉ số đã có cho nhiệm vụ này.}


\clearpage
\section{Kiến trúc mô hình BR-Dummy: Geo-I, mạng đường và các cơ chế}\label{sec:2}
\textbf{Nền tảng là Geo-I:} từ GPS thật, mô hình lấy ngẫu nhiên một vị trí đại diện đã được bảo vệ, gọi là \textbf{neo}. Lịch sử neo, mạng đường và POI công khai được dùng để sinh K vị trí/quỹ đạo giả. Phản hồi máy chủ được cache theo epoch; GPS thật xếp hạng top-5 riêng tại thiết bị. Lớp xử lý biên (boundary) bổ sung bảo vệ điểm đầu/cuối.

\begin{center}
\includegraphics[width=\linewidth]{figures/architecture_model.pdf}
\end{center}
Đọc từ trái sang phải: input $\to$ layer bảo vệ Geo-I $\to$ layer ngữ cảnh $\to$ layer sinh dummy $\to$ output. Khung xanh bao toàn bộ mô hình của ta, gồm ba layer lõi và layer 4 boundary S9/S10 ở hàng dưới; input và bên nhận ở ngoài khung. GeoI-Paced/GeoI-Slack dùng pacing; bản Slack thêm slack. Cache cho output top-5 riêng tại thiết bị.

\begingroup\small
\begin{longtable}{@{}P{4.105cm}P{12.814cm}@{}}
\toprule
\textbf{Component} & \textbf{Cơ chế và vai trò đối với scenario}\\\midrule
\endfirsthead
\toprule
\textbf{Component} & \textbf{Cơ chế và vai trò đối với scenario}\\\midrule
\endhead
\bottomrule\endfoot
Geo-I / neo REM & Lấy neo bằng cơ chế mũ REM: vị trí càng xa GPS theo khoảng cách Euclid, xác suất được chọn càng thấp. Tập ứng viên đường được xác định công khai; không cắt ứng viên theo bán kính riêng quanh GPS. Đây là nền bảo vệ S1.\\
\addlinespace[3pt]
Tái dùng neo + ngân sách & Kiểm tra khoảng cách có nhiễu để giữ neo cũ hoặc tạo neo mới; cả hai bước đều tiêu ngân sách riêng tư. Tái dùng neo giúp giảm số mẫu độc lập khi dừng (S2). Ngân sách giới hạn rò rỉ tích lũy qua chuỗi (S3); pacing giãn các lần đọc GPS.\\
\addlinespace[3pt]
Belief từ lịch sử neo & Belief là phân bố ước lượng xe có thể đang ở đâu, được cập nhật từ neo đã bảo vệ. Switching dừng/đi là biến thể trước, không dùng để tạo các số hiện tại. Khối này hỗ trợ chọn dummy cho S2/S3, không dùng tốc độ thật hay đích tương lai.\\
\addlinespace[3pt]
Road-aware + POI-aware & Road-aware giới hạn dummy ở nơi có thể tới theo hướng làn và luật rẽ. POI-aware chọn K điểm để phủ nhiều POI bổ sung nhau. Exchange đổi ứng viên; slack cho phép độ lệch khi chọn, tùy biến thể. Hỗ trợ chuỗi hợp lý (S3) và dịch vụ; sức chống suy luận vẫn cần đo.\\
\addlinespace[3pt]
S9: bỏ đầu trước lõi & Không đưa GPS của h giây đầu vào lõi. Vì vậy, những điểm đầu bị che không đi vào trạng thái neo rồi ảnh hưởng các bản tin sau. Đối thủ vẫn có thể suy ngược từ tuyến còn thấy.\\
\addlinespace[3pt]
S10: buffer sau lõi & Giữ mỗi bản tin đã bảo vệ trong bộ đệm $\Delta$ giây rồi mới phát. Khi chuyến kết thúc, hủy bản tin chưa phát. Cách này che phần cuối mà không cần biết trước đích, nhưng làm tăng độ trễ và chưa loại hết suy luận từ phần còn thấy.\\
\end{longtable}
\endgroup
Belief, mạng đường và bộ chọn POI chỉ xử lý neo đã bảo vệ cùng dữ liệu công khai. Theo tính chất hậu xử lý, chúng giữ bảo đảm của neo khi tính đúng ngân sách qua các lần truy cập (composition). K dummy không tự tạo k-anonymity. GPS thật chỉ thêm ở bước chọn top-5 tại thiết bị sau phản hồi, không đưa vào bộ chọn dummy. Pacing chưa che giờ chuyến. Boundary đã kiểm tra mã/tích hợp; chưa có benchmark toàn mô hình với dịch vụ.


\clearpage
\subsection{Định lượng hai cấu hình Geo-I và đóng góp của mô hình}\label{sub:2.1}
Hai tên dễ nhớ dưới đây đều thuộc \textbf{BR-Dummy trên nền Geo-I}, dùng cùng ngân sách riêng tư. GeoI-Paced là mốc so sánh nội bộ; GeoI-Slack là cấu hình làm việc hiện tại. Chúng khác cách chọn dummy, không phải hai mức $\varepsilon$ ưu tiên privacy/utility khác nhau.

\begingroup\small
\begin{longtable}{@{}P{3.938cm}P{6.350cm}P{6.350cm}@{}}
\toprule
\textbf{Tham số / component} & \textbf{GeoI-Paced} & \textbf{GeoI-Slack}\\\midrule
\endfirsthead
\toprule
\textbf{Tham số / component} & \textbf{GeoI-Paced} & \textbf{GeoI-Slack}\\\midrule
\endhead
\bottomrule\endfoot
Mã thực nghiệm & response\_\allowbreak{}paced / epoch\_\allowbreak{}cache & response\_\allowbreak{}paced\_\allowbreak{}slack03 / epoch\_\allowbreak{}cache\\
\addlinespace[3pt]
Neo riêng tư & REM; $\varepsilon$\_\allowbreak{}test = $\varepsilon$\_\allowbreak{}release = 0,01 /m; θ = 200 m & Cùng tham số neo\\
\addlinespace[3pt]
Ngân sách phiên & 23 đơn vị $\times$ 0,01; cận phiên 0,23 /m & Cùng cận phiên 0,23 /m\\
\addlinespace[3pt]
Pacing đọc GPS & Cách nhau ít nhất 60 giây & Cùng nhịp đọc\\
\addlinespace[3pt]
Road / belief / POI & Belief từ neo; làn có hướng, luật rẽ; chọn tập phủ POI & Cùng nền; thêm bước di chuyển có slack\\
\addlinespace[3pt]
Slack trong bộ chọn & Không nới objective & Cho phép mất tối đa 0,03 objective để tiến về mục tiêu; không phải giảm Recall 3\%\\
\addlinespace[3pt]
Đầu ra / dịch vụ & K=5 tọa độ; 6 loại POI; server trả top-10 mỗi loại/điểm & Cùng K=5, L=10\\
\addlinespace[3pt]
Cache và kết quả riêng & Hợp phản hồi cùng epoch 60 giây; GPS thật xếp hạng top-5 tại thiết bị & Cùng cache; không thay query\\
\end{longtable}
\endgroup
Lần đầu tạo neo chi 1 đơn vị; trước bước sau dự trữ 2 đơn vị. Test tái dùng chi 1, test + tạo neo mới chi 2. Giữa hai lần đọc hoặc khi hết ngân sách, mô hình chỉ dùng trạng thái đã bảo vệ. Nhịp đọc GPS, nhịp query dịch vụ và epoch hiệu lực phản hồi là ba khái niệm riêng; pacing không che giờ bắt đầu/kết thúc chuyến.

Trong phân tích lý tưởng với đồng hồ cố định, cận tỷ số xác suất của transcript là exp(0,23 $\times$ D$\infty$), với D$\infty$ là độ lệch GPS lớn nhất giữa hai chuỗi. Đây là cận theo khoảng cách, không phải xác suất “an toàn 23\%”. Hậu xử lý từ neo giữ cận này; mã số thực chưa có chứng nhận pure-DP chính xác.

\subsection{Đóng góp cần trình bày như thế nào?}\label{sub:2.2}
\textbf{Đề xuất:} từ lịch sử Geo-I đã bảo vệ, sinh tập truy vấn đi được trên mạng làn và phủ POI hữu ích trong ngân sách hữu hạn. Giá trị nằm ở sự phối hợp cơ chế theo nhiệm vụ dịch vụ: giảm suy luận gần, giữ khả năng tìm đúng địa điểm và kiểm soát tải truy vấn. Geo-I, tái dùng neo và dummy đã có tiền lệ; paper EV [\hyperlink{ref-R12}{R12}] đã kết hợp AGeoI + dummy, nên riêng tổ hợp hai từ này chưa là điểm mới.

\textbf{So với đối chứng:} DLS chú trọng độ khó phân biệt giữa điểm thật/dummy; các hướng TransProtect/Semantic chú trọng tính hợp lý của vị trí/chuỗi. Ta bổ sung việc chọn cả tập truy vấn theo hợp POI từ belief đã bảo vệ, thay vì dùng độ giống quỹ đạo làm đại diện cho utility. Road-aware tạo tính khả thi; privacy vẫn phải kiểm tra bằng attacker. Với RDG và Fake-query, đây mới là khác biệt thiết kế: chưa có benchmark Geo-I cùng protocol để chứng minh ưu thế số điểm.

\key{Lập luận đóng góp phải có cả cơ chế và bằng chứng: so trực tiếp bản Geo-I v2 với adapter để thấy đánh đổi; dùng phép bỏ/thêm component trên bản mới để kiểm tra lợi ích dịch vụ. Không gộp hai phiên bản thành một kết quả.}


\clearpage
\section{Benchmark Geo-I: so trực tiếp với đối chứng}\label{sec:3}
\textbf{Phép thử v2 đã được kiểm chứng:} K=5, ba seed 81/82/83, 12 chuyến test riêng biệt; mỗi chuyến có năm cửa sổ S1/S2/S3/S9/S10, nên 60 cửa sổ không phải 60 mẫu độc lập. Bản BR tái dùng dùng nền Geo-I với cận phiên 0,24 /m trên đồ thị nút giao. Đây là phiên bản trước mạng làn/pacing/slack/cache hiện tại, không có phân ca A/B/C.

* DLS là bản thích nghi lên đồ thị; TransProtect và Semantic là adapter rút gọn trong benchmark v2, không tái lập Transformer/LSTM của paper. Không thay chúng bằng adapter mới rồi giữ số cũ. RDG/Fake-query chưa có kết quả so trực tiếp với Geo-I trong phép thử này. AnotherMe chỉ áp dụng S3: 11/12 thành công, không cùng mẫu số nên báo riêng, không đưa vào bảng chung.

\begingroup\small
\begin{longtable}{@{}P{3.364cm}P{2.486cm}P{2.486cm}P{2.486cm}P{2.486cm}P{2.486cm}@{}}
\toprule
\textbf{Phương pháp} & \textbf{S1} & \textbf{S2} & \textbf{S3} & \textbf{S9} & \textbf{S10}\\\midrule
\endfirsthead
\toprule
\textbf{Phương pháp} & \textbf{S1} & \textbf{S2} & \textbf{S3} & \textbf{S9} & \textbf{S10}\\\midrule
\endhead
\bottomrule\endfoot
GPS thật & 100,00\% / 13 & 100,00\% / 39 & 96,43\% / 23 & 41,67\% / 120 & 8,33\% / 630\\
\addlinespace[3pt]
DLS* & 58,33\% / 1.318 & 100,00\% / 39 & 31,21\% / 1.688 & 8,33\% / 1.582 & 0,00\% / 2.029\\
\addlinespace[3pt]
TransProtect* & 83,33\% / 84 & 66,67\% / 88 & 69,28\% / 99 & 25,00\% / 184 & 8,33\% / 680\\
\addlinespace[3pt]
Semantic* & 100,00\% / 11 & 100,00\% / 42 & 40,58\% / 257 & 33,33\% / 132 & 16,67\% / 731\\
\addlinespace[3pt]
Geo-I + dummy & 58,33\% / 94 & 36,11\% / 105 & 52,65\% / 144 & 8,33\% / 307 & 0,00\% / 636\\
\addlinespace[3pt]
Geo-I + dummy đường & 16,67\% / 237 & 53,70\% / 121 & 3,57\% / 1.909 & 41,67\% / 217 & 0,00\% / 2.682\\
\addlinespace[3pt]
\textbf{BR tái dùng (Geo-I, v2)} & 33,33\% / 299 & 22,22\% / 247 & 7,44\% / 776 & 0,00\% / 316 & 0,00\% / 2.315\\
\end{longtable}
\endgroup
Mỗi ô: \textbf{Hit100 (\%) $\downarrow$ / MAE (m) $\uparrow$}. MAE và Hit chọn attacker riêng từ tập chọn đối thủ; không phải kết quả của một đối thủ duy nhất tối ưu cả hai.

\begingroup\small
\begin{longtable}{@{}P{3.364cm}P{2.486cm}P{2.486cm}P{2.486cm}P{2.486cm}P{2.486cm}@{}}
\toprule
\textbf{Recall@5 $\uparrow$} & \textbf{S1} & \textbf{S2} & \textbf{S3} & \textbf{S9} & \textbf{S10}\\\midrule
\endfirsthead
\toprule
\textbf{Recall@5 $\uparrow$} & \textbf{S1} & \textbf{S2} & \textbf{S3} & \textbf{S9} & \textbf{S10}\\\midrule
\endhead
\bottomrule\endfoot
GPS thật & 100,00\% & 100,00\% & 100,00\% & 100,00\% & 100,00\%\\
\addlinespace[3pt]
DLS* & 100,00\% & 100,00\% & 100,00\% & 100,00\% & 100,00\%\\
\addlinespace[3pt]
TransProtect* & 97,78\% & 97,19\% & 97,93\% & 97,27\% & 98,52\%\\
\addlinespace[3pt]
Semantic* & 100,00\% & 100,00\% & 100,00\% & 100,00\% & 100,00\%\\
\addlinespace[3pt]
Geo-I + dummy & 99,17\% & 98,52\% & 97,73\% & 96,85\% & 97,06\%\\
\addlinespace[3pt]
Geo-I + dummy đường & 98,06\% & 97,44\% & 81,96\% & 96,99\% & 79,75\%\\
\addlinespace[3pt]
\textbf{BR tái dùng (Geo-I, v2)} & 95,83\% & 97,01\% & 87,67\% & 96,60\% & 88,33\%\\
\end{longtable}
\endgroup
\textbf{Bằng chứng hỗ trợ:} ở S2, BR tái dùng có Hit100 \textbf{22,22\%}, thấp hơn DLS/Semantic (100\%) và TransProtect (66,67\%), với Recall 97,01\%. Ở S3, Hit100 \textbf{7,44\%}, thấp hơn ba adapter (31,21–69,28\%) và Geo-I + dummy (52,65\%). Kết quả hỗ trợ thiết kế phối hợp trạng thái đã bảo vệ theo thời gian; chưa cô lập riêng tác động tái dùng vì tham số và cách sinh dummy cũng khác.

\textbf{Đánh đổi phải nói rõ:} BR không thắng mọi metric. Geo-I + dummy đường có Hit100 S3 thấp hơn (3,57\%) nhưng Recall thấp hơn (81,96\% so với 87,67\%); cả hai chưa đạt 90\%. DLS có MAE lớn hơn ở S1/S3/S9 dù Hit100 có thể cao hơn. BR cũng chưa đạt Recall 90\% ở S10 (88,33\%); S10 Hit100=0 của nhiều phương pháp nên không đủ phân biệt. S9/S10 chỉ suy tọa độ đầu/cuối SUMO, chưa chứng minh che nhà hay nơi làm việc.

\key{Kết luận v2: BR trên nền Geo-I giảm tỷ lệ suy đúng gần ở một số nhiệm vụ quan trọng, có đánh đổi utility. Đây là bằng chứng đóng góp của bản thích nghi trong benchmark, chưa chứng minh vượt các paper nguyên bản hoặc ưu thế ở cùng ngân sách.}


\clearpage
\subsection{Bản Geo-I hiện tại: 14 ca và kiểm tra từng component}\label{sub:3.1}
\textbf{GeoI-Slack:} 12 nhóm phát triển, 165 bản ghi từ 94 chuyến nguồn; S10 chỉ A/B. K=5, L=10, cận phiên 0,23 /m. Recall dưới đây dùng POI khả dụng mô phỏng ở p=0,8, ba thế giới khả dụng và hai lần chạy cơ chế. Privacy kế thừa từ đúng tọa độ công bố đã chấm bằng tập attacker hữu hạn: cache và trạng thái server độc lập với GPS không đổi query. Đây là dữ liệu phát triển, không phải xác nhận trên nhóm mới.

\begingroup\small
\begin{longtable}{@{}P{1.589cm}P{2.980cm}P{3.477cm}P{2.980cm}P{3.179cm}P{1.589cm}@{}}
\toprule
\textbf{Ca} & \textbf{GPS thật: Hit100} & \textbf{GeoI-Slack: Hit100 $\downarrow$} & \textbf{MAE (m) $\uparrow$} & \textbf{Recall@5 $\uparrow$} & \textbf{Nhóm}\\\midrule
\endfirsthead
\toprule
\textbf{Ca} & \textbf{GPS thật: Hit100} & \textbf{GeoI-Slack: Hit100 $\downarrow$} & \textbf{MAE (m) $\uparrow$} & \textbf{Recall@5 $\uparrow$} & \textbf{Nhóm}\\\midrule
\endhead
\bottomrule\endfoot
S1.A & 100,00\% & 8,33\% & 852 & 97,64\% & 12\\
\addlinespace[3pt]
S1.B & 100,00\% & 8,33\% & 832 & 95,90\% & 12\\
\addlinespace[3pt]
S1.C & 100,00\% & 0,00\% & 1.674 & 80,18\% & 12\\
\addlinespace[3pt]
S2.A & 100,00\% & 0,00\% & 910 & 99,86\% & 12\\
\addlinespace[3pt]
S2.B & 100,00\% & 4,17\% & 715 & 99,49\% & 12\\
\addlinespace[3pt]
S2.C & 100,00\% & 0,00\% & 997 & 94,59\% & 12\\
\addlinespace[3pt]
S3.A & 100,00\% & 3,09\% & 906 & 95,63\% & 12\\
\addlinespace[3pt]
S3.B & 100,00\% & 3,37\% & 873 & 95,97\% & 11\\
\addlinespace[3pt]
S3.C & 100,00\% & 2,19\% & 934 & 95,26\% & 12\\
\addlinespace[3pt]
S9.A & 16,67\% & 0,00\% & 607 & 95,69\% & 12\\
\addlinespace[3pt]
S9.B & 75,00\% & 2,50\% & 701 & 95,30\% & 10\\
\addlinespace[3pt]
S9.C & 0,00\% & 0,00\% & 514 & 94,95\% & 12\\
\addlinespace[3pt]
S10.A & 33,33\% & 0,00\% & 1.345 & 97,63\% & 12\\
\addlinespace[3pt]
S10.B & 16,67\% & 0,00\% & 1.200 & 96,42\% & 12\\
\end{longtable}
\endgroup
\begingroup\small
\begin{longtable}{@{}P{4.839cm}P{3.581cm}P{3.194cm}P{4.742cm}@{}}
\toprule
\textbf{Cấu hình} & \textbf{Recall@5 $\uparrow$} & \textbf{Ca $\ge$90\%} & \textbf{Byte / sự kiện $\downarrow$}\\\midrule
\endfirsthead
\toprule
\textbf{Cấu hình} & \textbf{Recall@5 $\uparrow$} & \textbf{Ca $\ge$90\%} & \textbf{Byte / sự kiện $\downarrow$}\\\midrule
\endhead
\bottomrule\endfoot
GeoI-Paced & 94,38\% & 13/14 & 993,7\\
\addlinespace[3pt]
\textbf{GeoI-Slack} & 95,44\% & 13/14 & 994,1\\
\addlinespace[3pt]
GeoI-Slack, bỏ cache & 95,19\% & 13/14 & 994,1\\
\addlinespace[3pt]
GPS thật & 100,00\% & 14/14 & 214,5\\
\end{longtable}
\endgroup
\textbf{Đóng góp có thể định lượng:} cache tăng Recall 0,24 điểm \%, CI 95\% [0,14; 0,34], không thêm query/byte trong schema này. Thêm slack tăng trung bình 1,06 điểm \%, CI [-0,13; 2,23] chứa 0: hướng cải thiện có triển vọng, chưa đủ kết luận chắc chắn. CI lấy mẫu lại ghép cặp 12 nhóm 10.000 lần, chưa hiệu chỉnh nhiều so sánh.

\textbf{Giới hạn còn thấy:} S1.C chỉ đạt 80,18\% Recall, nên cấu hình đạt 13/14 ca $\ge$90\%, không phải đạt mọi ca. S9.C có Hit100=0 cả với GPS thật, nên riêng chỉ số này chưa chứng minh cải thiện; Geo-I vẫn có Hit500=58,33\%: điểm đầu vẫn có thể bị khoanh vùng. Những số S9/S10 là của lõi Geo-I qua cửa sổ quan sát; không gán cho BR-Boundary, vì wrapper cắt đầu/buffer đuôi chưa được benchmark cùng dịch vụ này.

Recall gộp đều ca trong scenario rồi đều năm scenario. Byte gộp theo sự kiện, tính toàn bộ traffic của 94 chuyến được giữ, gồm yêu cầu + phản hồi JSON; chưa gồm HTTP/TLS hoặc latency. S10.B ở đây giữ nguồn cũ S10.C. Các số tổng/CI được tính lại đúng 14 ca, không chạy lại mô hình và không trộn với kết quả v2.


\clearpage
\section{Phụ lục: đọc chi tiết 14 sample A/B/C}\label{sec:4}
A/B/C là các điều kiện cùng nhiệm vụ, không phải thứ tự độ khó. Mỗi ca chọn một bản ghi thật từ bộ minh họa 12 nhóm/264 chuyến; số bản ghi không phải số chuyến độc lập. \textbf{S10 chỉ có A/B} trong tài liệu; B ánh xạ mã nguồn cũ S10.C để giữ truy nguyên, không có ca C độc lập trong phạm vi hiện tại.

Chấm màu = mẫu trước bảo vệ; nét đứt = đường thật để chấm; sao đỏ = đáp án; trục thời gian tách các mẫu chồng nhau. Đối thủ chỉ nhận dữ liệu đã bảo vệ và thông tin phụ trợ được phép, không nhận các tọa độ/nhãn thật này. Bản đồ minh họa dùng nền SUMO với cùng mã kiểm tra tệp OSM; thiếu mạng gốc để xác minh trùng toàn bộ hình học. © OpenStreetMap contributors.

\subsection{S1 — suy ra vị trí tại một lần gửi}\label{sub:4.1}
Đối thủ cần suy ra tọa độ tại một lần gửi. A/B thay ràng buộc mạng đường; C thêm ngữ cảnh POI hiếm.

\noindent\begin{minipage}[t]{0.48\linewidth}\vspace{0pt}\includegraphics[width=\linewidth]{figures/case_S1_A.pdf}\end{minipage}\hfill\begin{minipage}[t]{0.49\linewidth}\vspace{0pt}\RaggedRight \textbf{S1.A / 12 bản ghi.} Một bản tin tại cạnh có $\ge$2 hướng đi tiếp hợp lệ: kiểm tra khi mạng đường còn nhiều khả năng.\par\smallskip \textbf{Bản ghi minh họa v2-r301-000.} u301\_\allowbreak{}00: 1 mẫu, chỉ số GPS FCD[333…333]\par\smallskip\textbf{Đọc sample.} Một lần lấy mẫu ở cạnh có 4 hướng đi tiếp. Sao đỏ trùng vị trí đầu vào; đối thủ phải chọn vị trí từ đầu ra đã bảo vệ, không được thấy sẵn sao này.\end{minipage}\par\vspace{3pt}
\noindent\begin{minipage}[t]{0.48\linewidth}\vspace{0pt}\includegraphics[width=\linewidth]{figures/case_S1_B.pdf}\end{minipage}\hfill\begin{minipage}[t]{0.49\linewidth}\vspace{0pt}\RaggedRight \textbf{S1.B / 12 bản ghi.} Một bản tin tại cạnh chỉ có 1 hướng đi tiếp: bản đồ tự loại bớt khả năng.\par\smallskip \textbf{Bản ghi minh họa v2-r301-001.} u301\_\allowbreak{}00: 1 mẫu, chỉ số GPS FCD[241…241]\par\smallskip\textbf{Đọc sample.} Vẫn chỉ một lần lấy mẫu, nhưng cạnh chỉ có một hướng đi tiếp. Bản đồ có thể loại bớt khả năng so với A; cần đo lợi thế này bằng đối thủ thực tế.\end{minipage}\par\vspace{3pt}
\noindent\begin{minipage}[t]{0.48\linewidth}\vspace{0pt}\includegraphics[width=\linewidth]{figures/case_S1_C.pdf}\end{minipage}\hfill\begin{minipage}[t]{0.49\linewidth}\vspace{0pt}\RaggedRight \textbf{S1.C / 12 bản ghi.} Cách POI $\le$150 m; loại POI chiếm $\le$5\% danh mục công khai. Hiếm theo loại địa điểm, không phải theo tần suất ghé thăm.\par\smallskip \textbf{Bản ghi minh họa v2-r301-029.} u301\_\allowbreak{}13: 1 mẫu, chỉ số GPS FCD[695…695]\par\smallskip\textbf{Đọc sample.} Điểm lấy mẫu cách POI pharmacy 12,03 m; loại này chiếm 3,11\% danh mục. Ngữ cảnh địa điểm hiếm là dấu hiệu bổ sung, không phải chuỗi di chuyển.\end{minipage}\par\vspace{3pt}

\clearpage
\subsection{S2 — suy ra nơi dừng qua nhiều lần gửi}\label{sub:4.2}
Đối thủ cần suy ra tọa độ nơi dừng. A/B thay thời lượng và nhịp quan sát; C kiểm tra rời đi rồi quay lại.

\noindent\begin{minipage}[t]{0.48\linewidth}\vspace{0pt}\includegraphics[width=\linewidth]{figures/case_S2_A.pdf}\end{minipage}\hfill\begin{minipage}[t]{0.49\linewidth}\vspace{0pt}\RaggedRight \textbf{S2.A / 12 bản ghi.} Dừng có chủ đích $\ge$20 giây; lấy truy vấn mỗi 5 giây.\par\smallskip \textbf{Bản ghi minh họa v2-r301-002.} u301\_\allowbreak{}05: 6 mẫu, chỉ số GPS FCD[54…79]\par\smallskip\textbf{Đọc sample.} 6 mẫu cùng một tọa độ, cách nhau 5 giây; đoạn dừng thật dài 29 giây. Sáu vạch thời gian bị chồng thành một chấm trên map.\end{minipage}\par\vspace{3pt}
\noindent\begin{minipage}[t]{0.48\linewidth}\vspace{0pt}\includegraphics[width=\linewidth]{figures/case_S2_B.pdf}\end{minipage}\hfill\begin{minipage}[t]{0.49\linewidth}\vspace{0pt}\RaggedRight \textbf{S2.B / 12 bản ghi.} Dừng $\ge$120 giây; lấy truy vấn mỗi 20 giây.\par\smallskip \textbf{Bản ghi minh họa v2-r301-003.} u301\_\allowbreak{}06: 9 mẫu, chỉ số GPS FCD[59…219]\par\smallskip\textbf{Đọc sample.} 9 mẫu cùng một tọa độ, cách nhau 20 giây; đoạn dừng dài 179 giây. Đối thủ có cửa sổ dài hơn để kết hợp bằng chứng, không chỉ một bản tin.\end{minipage}\par\vspace{3pt}
\noindent\begin{minipage}[t]{0.48\linewidth}\vspace{0pt}\includegraphics[width=\linewidth]{figures/case_S2_C.pdf}\end{minipage}\hfill\begin{minipage}[t]{0.49\linewidth}\vspace{0pt}\RaggedRight \textbf{S2.C / 12 bản ghi.} Hai lần dừng trên cùng làn, mỗi lần $\ge$20 giây; ở giữa thực sự di chuyển.\par\smallskip \textbf{Bản ghi minh họa v2-r301-004.} u301\_\allowbreak{}07: 18 mẫu, chỉ số GPS FCD[214…1193]\par\smallskip\textbf{Đọc sample.} 18 mẫu chia hai đợt, mỗi đợt 9 mẫu. Hai lần dừng dài 44 giây tại cùng tọa độ; khoảng trống giữa hai cụm thời gian là lúc xe rời đi rồi quay lại.\end{minipage}\par\vspace{3pt}

\clearpage
\subsection{S3 — tái dựng đoạn đường đã đi}\label{sub:4.3}
Đối thủ cần tái dựng các vị trí của đoạn đã đi. A là chuỗi đều; B ít nhánh; C quan sát thưa.

\noindent\begin{minipage}[t]{0.48\linewidth}\vspace{0pt}\includegraphics[width=\linewidth]{figures/case_S3_A.pdf}\end{minipage}\hfill\begin{minipage}[t]{0.49\linewidth}\vspace{0pt}\RaggedRight \textbf{S3.A / 12 bản ghi.} Đoạn di chuyển $\ge$60 giây qua $\ge$3 cạnh; gửi mỗi 20 giây.\par\smallskip \textbf{Bản ghi minh họa v2-r301-005.} u301\_\allowbreak{}00: 12 mẫu, chỉ số GPS FCD[361…581]\par\smallskip\textbf{Đọc sample.} 12 mẫu của u301\_\allowbreak{}00 cách nhau 20 giây. Chuỗi điểm cung cấp nhiều ràng buộc để nối tuyến; nét đứt là đường thật chỉ dùng chấm.\end{minipage}\par\vspace{3pt}
\noindent\begin{minipage}[t]{0.48\linewidth}\vspace{0pt}\includegraphics[width=\linewidth]{figures/case_S3_B.pdf}\end{minipage}\hfill\begin{minipage}[t]{0.49\linewidth}\vspace{0pt}\RaggedRight \textbf{S3.B / 11 bản ghi.} Ít nhất một nửa số cạnh được lấy mẫu chỉ có một hướng đi tiếp: hành lang ít lựa chọn.\par\smallskip \textbf{Bản ghi minh họa v2-r301-007.} u301\_\allowbreak{}00: 9 mẫu, chỉ số GPS FCD[91…251]\par\smallskip\textbf{Đọc sample.} 9 mẫu; 55,6\% cạnh được lấy mẫu chỉ có một hướng đi tiếp. Đường ít nhánh có thể khiến chuỗi dễ suy hơn dù từng vị trí được làm mờ.\end{minipage}\par\vspace{3pt}
\noindent\begin{minipage}[t]{0.48\linewidth}\vspace{0pt}\includegraphics[width=\linewidth]{figures/case_S3_C.pdf}\end{minipage}\hfill\begin{minipage}[t]{0.49\linewidth}\vspace{0pt}\RaggedRight \textbf{S3.C / 12 bản ghi.} Lấy mẫu đoạn di chuyển thưa hơn, mỗi 60 giây.\par\smallskip \textbf{Bản ghi minh họa v2-r301-006.} u301\_\allowbreak{}00: 4 mẫu, chỉ số GPS FCD[361…541]\par\smallskip\textbf{Đọc sample.} 4 mẫu của cùng đoạn trong A, cách nhau 60 giây. Khoảng trống lớn hơn nhưng đối thủ vẫn có thể dùng mạng đường để nối các khả năng.\end{minipage}\par\vspace{3pt}

\clearpage
\subsection{S9 — suy ra điểm xuất phát bị che}\label{sub:4.4}
Đối thủ cần suy điểm đầu bị giấu. A suy ngược một chuyến; B phân biệt hai nguồn nhập tuyến; C kết hợp chuyến lặp.

\noindent\begin{minipage}[t]{0.48\linewidth}\vspace{0pt}\includegraphics[width=\linewidth]{figures/case_S9_A.pdf}\end{minipage}\hfill\begin{minipage}[t]{0.49\linewidth}\vspace{0pt}\RaggedRight \textbf{S9.A / 12 bản ghi.} Cạnh xuất phát chỉ một lối ra. Che 60 giây đầu, lấy phần còn lại mỗi 20 giây.\par\smallskip \textbf{Bản ghi minh họa v2-r301-025.} u301\_\allowbreak{}00: 27 mẫu, chỉ số GPS FCD[60…580]\par\smallskip\textbf{Đọc sample.} Sao là FCD[0]; cửa sổ bắt đầu ở FCD[60] sau khi che 60 giây. Chỉ một lối ra có thể giúp suy ngược nguồn xuất phát từ phần còn công bố.\end{minipage}\par\vspace{3pt}
\noindent\begin{minipage}[t]{0.48\linewidth}\vspace{0pt}\includegraphics[width=\linewidth]{figures/case_S9_B.pdf}\end{minipage}\hfill\begin{minipage}[t]{0.49\linewidth}\vspace{0pt}\RaggedRight \textbf{S9.B / 11 bản ghi.} Hai chuyến khác điểm đầu nhưng nhập một đoạn chung; chỉ cấp từ chỗ nhập tuyến.\par\smallskip \textbf{Bản ghi minh họa v2-r301-026.} u301\_\allowbreak{}00: 28 mẫu, chỉ số GPS FCD[34…574]; u301\_\allowbreak{}04: 28 mẫu, chỉ số GPS FCD[19…559]\par\smallskip\textbf{Đọc sample.} Hai điểm đầu cách 281,76 m rồi nhập tuyến. Chỉ cấp dữ liệu sau nhập; hai sao là hai đáp án bị giấu mà đối thủ phải phân biệt.\end{minipage}\par\vspace{3pt}
\noindent\begin{minipage}[t]{0.48\linewidth}\vspace{0pt}\includegraphics[width=\linewidth]{figures/case_S9_C.pdf}\end{minipage}\hfill\begin{minipage}[t]{0.49\linewidth}\vspace{0pt}\RaggedRight \textbf{S9.C / 12 bản ghi.} Nhiều chuyến có cùng điểm xuất phát dự kiến; che 60 giây đầu từng chuyến.\par\smallskip \textbf{Bản ghi minh họa v2-r301-027.} u301\_\allowbreak{}00: 27 mẫu, chỉ số GPS FCD[60…580]; u301\_\allowbreak{}09: 27 mẫu, chỉ số GPS FCD[60…580]\par\smallskip\textbf{Đọc sample.} u301\_\allowbreak{}00/u301\_\allowbreak{}09 lặp cùng điểm đầu dự kiến; mỗi phiên che 60 giây đầu. Nhiều cửa sổ có thể bổ sung bằng chứng về cùng nguồn xuất phát.\end{minipage}\par\vspace{3pt}

\clearpage
\subsection{S10 — suy ra điểm kết thúc bị che}\label{sub:4.5}
Đối thủ cần suy điểm cuối của chuyến đã hoàn tất. A dùng một chuyến; B kết hợp nhiều chuyến cùng đích. Đây khác dự đoán đích tương lai ở S6.

\noindent\begin{minipage}[t]{0.58\linewidth}\vspace{0pt}\includegraphics[width=\linewidth]{figures/case_S10_A.pdf}\end{minipage}\hfill\begin{minipage}[t]{0.39\linewidth}\vspace{0pt}\RaggedRight \textbf{S10.A / 11 bản ghi.} Cạnh kết thúc chỉ một lối vào. Che 60 giây cuối của chuyến đã hoàn tất.\par\smallskip \textbf{Bản ghi minh họa v2-r302-025.} u302\_\allowbreak{}08: 38 mẫu, chỉ số GPS FCD[0…740]\par\smallskip\textbf{Đọc sample.} Sao là FCD cuối [817] của u302\_\allowbreak{}08; che 60 giây cuối. Đường chỉ một lối vào có thể giúp suy đoạn kết thúc từ phần trước còn thấy.\end{minipage}\par\vspace{3pt}
\noindent\begin{minipage}[t]{0.58\linewidth}\vspace{0pt}\includegraphics[width=\linewidth]{figures/case_S10_B.pdf}\end{minipage}\hfill\begin{minipage}[t]{0.39\linewidth}\vspace{0pt}\RaggedRight \textbf{S10.B / 12 bản ghi.} Nhiều chuyến cùng điểm kết thúc dự kiến; che 60 giây cuối từng chuyến.\par\smallskip \textbf{Bản ghi minh họa v2-r301-028.} u301\_\allowbreak{}00: 27 mẫu, chỉ số GPS FCD[0…520]; u301\_\allowbreak{}09: 27 mẫu, chỉ số GPS FCD[0…520]\par\smallskip\textbf{Đọc sample.} u301\_\allowbreak{}00/u301\_\allowbreak{}09 lặp cùng đích dự kiến và cùng che 60 giây cuối. Đối thủ có thể kết hợp nhiều phiên để khoanh đích, chưa đủ gọi đó là nơi làm việc.\end{minipage}\par\vspace{3pt}
S9/S10 chấm tọa độ đầu/cuối; chưa gán nhãn nhà/nơi làm việc. Xem bản đồ phóng to của 29 ca của khung đầy đủ ở sample\_\allowbreak{}maps.html; tọa độ và bản ghi ở data\_\allowbreak{}samples.json, số đếm ở scenario\_\allowbreak{}inventory.csv.


\clearpage
\section{Nguồn và khả năng truy nguyên}\label{sec:5}
Kết quả so trực tiếp: artifacts/benchmarks/paper\_\allowbreak{}benchmark/results.json và docs/reviews/verification\_\allowbreak{}paper\_\allowbreak{}benchmark\_\allowbreak{}v2.md. Bản mới: iteration18\_\allowbreak{}expanded\_\allowbreak{}attacks.json, iteration28\_\allowbreak{}live\_\allowbreak{}service.json và iteration28\_\allowbreak{}verification.json trong artifacts/benchmarks/research\_\allowbreak{}loop/. geoi\_\allowbreak{}evidence.py kiểm tra hash và gộp lại đúng 14 ca; method\_\allowbreak{}evidence.json lưu nguồn, số đo và bảng đã in. preparation\_\allowbreak{}evidence.json truy nguyên guide. Không chạy thêm mô hình; hai phiên bản được báo riêng.

Ngày 03/10/2026 đưa Geo-I về phương pháp chính; giữ nguyên thực nghiệm gốc và báo riêng hai phiên bản. S10.B là nhãn cho mã nguồn S10.C cũ. Cấu hình, benchmark và phụ lục được dùng chung giữa report/guide; cách dựng và tệp tác giả ở README.md.

\par\hypertarget{ref-R1}{}\textbf{[R1]} Yadav et al. (2024). Protecting Vehicle Location Privacy with Contextually-Driven Synthetic Location Generation (TransProtect). SIGSPATIAL. \href{https://arxiv.org/html/2409.09495v1}{Nguồn gốc}. Toàn văn tác giả; §5.1.4, Eq. 13.
\par\hypertarget{ref-R2}{}\textbf{[R2]} Liu, Peng, Zhou (2026). A Dummy-Based Location Privacy Protection Scheme with Semantic Correlation of Moving Paths. JKSUCIS 38:478. \href{https://link.springer.com/article/10.1007/s44443-026-00899-w}{Nguồn gốc}. Toàn văn NXB; §6.3–6.5.
\par\hypertarget{ref-R3}{}\textbf{[R3]} Liu, Hu, Zhou (2026). Location Privacy Protection in Continuous LBSs: Enhancing Anonymity via Fake Queries. JKSUCIS 38:53. \href{https://link.springer.com/article/10.1007/s44443-025-00438-z}{Nguồn gốc}. Toàn văn NXB; §6.3–6.6.
\par\hypertarget{ref-R4}{}\textbf{[R4]} Road Network-Aware Personalized Trajectory Protection with Differential Privacy under Spatiotemporal Correlations (2025), arXiv:2511.21020v1. \href{https://arxiv.org/html/2511.21020v1}{Nguồn gốc}. Bản tiền công bố; §VI, Eq. 26–27.
\par\hypertarget{ref-R5}{}\textbf{[R5]} Wang et al. (2025). CPCROK: A Communication-Efficient and Privacy-Preserving Scheme for Low-Density Vehicular Ad Hoc Networks. Future Internet 17(4):165. \href{https://uhra.herts.ac.uk/id/eprint/25705/1/futureinternet-17-00165.pdf}{Nguồn gốc}. Toàn văn lưu tại trường tác giả; §5.4.
\par\hypertarget{ref-R6}{}\textbf{[R6]} Yue, Li, Liu, Li (2025). DP-FETC: A differentially private trajectory publishing method based on feature extraction and trajectory correlation. ERA 33(11):6631–6651. \href{https://www.aimspress.com/aimspress-data/era/2025/11/PDF/era-33-11-293.pdf}{Nguồn gốc}. Toàn văn NXB; §4.4, §5.2, Eq. 19–21.
\par\hypertarget{ref-R7}{}\textbf{[R7]} Gu et al. (2025). Research on Joint Protection of LBS Location and Query Privacy in Internet of Vehicles Based on an Improved PIR Algorithm. Concurrency and Computation: Practice and Experience. \href{https://onlinelibrary.wiley.com/doi/abs/10.1002/cpe.70447}{Nguồn gốc}. Chỉ xác minh tóm tắt NXB; chưa đủ định nghĩa các metric.
\par\hypertarget{ref-R8}{}\textbf{[R8]} Buchholz et al. (2024). SoK: Can Trajectory Generation Combine Privacy and Utility? PoPETs 2024(3):75–93. \href{https://petsymposium.org/popets/2024/popets-2024-0068.pdf}{Nguồn gốc}. Toàn văn; khảo sát và hướng dẫn đánh giá, không phải một cơ chế đối chứng.
\par\hypertarget{ref-R9}{}\textbf{[R9]} Farahnakiyan, Esmaeilyfard, Javidan (2025). A proactive privacy-preserving framework for mobile trajectory sharing (PRISM). JNCA 242:104271. \href{https://www.sciencedirect.com/science/article/pii/S1084804525001687}{Nguồn gốc}. Tóm tắt/đoạn công khai của NXB; chưa xác minh công thức LPI/TDD/DC.
\par\hypertarget{ref-R10}{}\textbf{[R10]} Ioannou, Aktypi, Athanasopoulos (2026). From Options to Action: Evaluating Adoption of Privacy Features in Fitness-Tracking Platforms. CHI 2026. \href{https://elathan.github.io/papers/chi26.pdf}{Nguồn gốc}. Toàn văn tác giả; bảng 3, §6.4. Đo sử dụng tính năng, không đo sức chống tấn công.
\par\hypertarget{ref-B1}{}\textbf{[B1]} Niu et al. (2014). Achieving k-Anonymity in Privacy-Aware Location-Based Services (DLS). INFOCOM. \href{https://doi.org/10.1109/INFOCOM.2014.6848002}{Nguồn gốc}. Đối chứng nền; định nghĩa entropy theo hồ sơ khảo sát đã có trong repo.
\par\hypertarget{ref-B2}{}\textbf{[B2]} Shaham et al. (2021). Privacy Preservation in Location-Based Services: A Novel Metric and Attack Model (RDG). IEEE TMC 20(10). \href{https://doi.org/10.1109/TMC.2020.2993599}{Nguồn gốc}. Đối chứng nền; entropy/chuyển tiếp/Viterbi theo hồ sơ đã có trong repo.
\par\hypertarget{ref-B3}{}\textbf{[B3]} Li et al. (2024 issue; online 11/09/2023). AnotherMe: A Location Privacy Protection System Based on Online Virtual Trajectory Generation. IEEE TDSC 21(4):2552–2567. \href{https://ieeexplore.ieee.org/document/10246991/}{Nguồn gốc}. Metadata/tóm tắt và mã nguồn tác giả; chưa kiểm chứng đầy đủ mẫu số/đơn vị phép đo.
\par\hypertarget{ref-B4}{}\textbf{[B4]} Brauer et al. (2022). My home is my secret: concealing sensitive locations by context-aware trajectory truncation. IJGIS 36(12):2496–2524. \href{https://doi.org/10.1080/13658816.2022.2081694}{Nguồn gốc}. Toàn văn NXB và chương tác giả 2023; S-TT, k-site-unidentifiability. Ngoại lệ nền trực tiếp cho S9/S10.
\par\hypertarget{ref-R11}{}\textbf{[R11]} Forsch et al. (2023). Efficient Mining of Volunteered Trajectory Datasets, §3.2. \href{https://link.springer.com/chapter/10.1007/978-3-031-35374-1_3}{Nguồn gốc}. Toàn văn tác giả, online 09/12/2023. Trình bày lại S-TT 2022, không phải một phương pháp mới độc lập.
\par\hypertarget{ref-B5}{}\textbf{[B5]} Dhondt et al. (2022). A Run a Day Won’t Keep the Hacker Away: Inference Attacks on Endpoint Privacy Zones in Fitness Tracking Social Networks. CCS. \href{https://people.cs.kuleuven.be/~stijn.volckaert/papers/2022_CCS_Fitness_Tracking.pdf}{Nguồn gốc}. Bản tác giả: tấn công EPZ và sáu countermeasures; không gắn nhãn một model bảo vệ đã thắng benchmark.
\par\hypertarget{ref-B6}{}\textbf{[B6]} Dehaene et al. (2023). Masking Location Streams in the Presence of Colluding Service Providers. EuroS\&P Workshops. \href{https://lirias.kuleuven.be/bitstream/20.500.12942/720391/2/Data_Stream_Obfuscator_camera_ready.pdf}{Nguồn gốc}. Bản tác giả: Data Stream Obfuscator, privacy zones/timeslots, mapping cache. Tháng 7/2023, ngoài cửa sổ ba năm.
\par\hypertarget{ref-R12}{}\textbf{[R12]} Atmaca et al. (2024). A Privacy-Preserving Querying Mechanism with High Utility for Electric Vehicles. IEEE OJVT. \href{https://wrap.warwick.ac.uk/id/eprint/183198/1/WRAP-privacy-preserving-querying-mechanism-high-utility-electric-vehicles-2024.pdf}{Nguồn gốc}. Bản tác giả: AGeoI + dummy cho truy vấn trạm sạc; không xác minh benchmark điểm đầu/cuối.
\end{document}
